\documentclass[11pt,a4paper]{article}
\usepackage[T1]{fontenc}
\usepackage{amsmath,amssymb,geometry,booktabs,array,parskip,xcolor,hyperref}
\geometry{margin=2.3cm}
\hypersetup{colorlinks=true,linkcolor=blue!60!black,urlcolor=blue!60!black}
\newcolumntype{L}[1]{>{\raggedright\arraybackslash}p{#1}}
\title{\textbf{What We Are Doing, and Why}\\\large A physics student's guide to the project}
\author{Companion to WorkPlan.tex}
\date{Verified results: 24 October 2026}
\begin{document}
\begin{center}
{\LARGE\bfseries What We Are Doing, and Why}\\[0.5em]
{\large A physics student's guide to the project}\\[0.4em]
{\small Companion to WorkPlan.tex; verified results: 24 October 2026}
\end{center}

\section*{The question we are trying to answer}

Suppose we hold two different stocks and want to sell them by a deadline.
Selling everything immediately can move the market price against us.
Selling slowly reduces that effect, but leaves us exposed to price changes
for longer. We want the selling schedule that best balances these costs.

This is our \textbf{test problem}. The research question is then:
\emph{how many qubits and quantum operations would a quantum computer need
to solve a reformulation of this problem?}
We already have a classical solution, so we can use it to check each later
step. We are estimating the quantum cost; we have not demonstrated a quantum
speedup.

The route is:
\begin{center}
trading cost $\longrightarrow$ optimal-control equation\\
$\longrightarrow$ linear diffusion equation $\longrightarrow$ quantum simulation\\
$\longrightarrow$ qubit and gate counts.
\end{center}

\section*{The quantities to keep in mind}

\begin{tabular}{@{}L{2.8cm}L{12cm}@{}}
\toprule
Symbol & Meaning \\
\midrule
$\mathbf x(t)=(x_1,x_2)^\top$ & Inventory: how much of each stock we still hold. \\
$\mathbf X=\mathbf x(0)$ & The inventory at the beginning. \\
$\mathbf v(t)=-\dot{\mathbf x}(t)$ & Selling rate. Positive means selling; negative means buying. \\
$T$ & The deadline. Exact liquidation means $\mathbf x(T)=0$. \\
$H$ & Matrix describing the cost of trading quickly. \\
$C$ & Covariance matrix: price uncertainty and how the two prices move together. \\
$\lambda$ & How strongly we penalize risk relative to trading cost. \\
\bottomrule
\end{tabular}

Bold letters are vectors. A transpose, $\top$, turns a column into a row,
so $\mathbf x^\top C\mathbf x$ is a single number, like the quadratic forms
used for kinetic or potential energy.

\section*{Step 1 --- Start with a list of trades}

Divide the available time into $N$ intervals of length $\tau=T/N$.
If $n_{jk}$ is the amount of stock $j$ sold in interval $k$, then
\[
 n_{jk}=x_{j,k-1}-x_{jk},\qquad
 U=\underbrace{\mathbb E[\text{cost}]}_{\text{average trading cost}}
   +\lambda\underbrace{\operatorname{Var}(\text{cost})}_{\text{risk}}.
\]
The expectation is an average over possible price changes. The variance
measures how uncertain the final cost is. This objective comes from
Almgren--Chriss~\cite{ac}; its expanded sum is in WorkPlan Stage~1.

\textbf{Where we are:} the discrete objective and its source have been checked.

\newpage
\section*{Step 2 --- Replace the trade list by a smooth path}

Let the time intervals become very small. A trade becomes a rate times a
time interval, $n_{jk}\approx v_j(t_k)\tau$, and the sum becomes an integral:
\begin{equation*}
 \boxed{J[\mathbf v]=\int_0^T
 \left(\mathbf v^\top H\mathbf v+\lambda\mathbf x^\top C\mathbf x\right)dt,}
 \qquad \dot{\mathbf x}=-\mathbf v.
\end{equation*}
This is a \emph{functional}: it assigns a number to an entire path, much like
an action in mechanics. Here we minimize a cost, not a mechanical action.

For diagonal $H=\operatorname{diag}(\eta_1,\eta_2)$, the first term is
$\eta_1v_1^2+\eta_2v_2^2$: trading twice as fast costs four times as much
per unit time. The second term penalizes inventory that remains exposed to
uncertain prices. For two assets,
\[
 \mathbf x^\top C\mathbf x=C_{11}x_1^2+2C_{12}x_1x_2+C_{22}x_2^2.
\]
The cross term couples the two assets. When their prices tend to move
together, holding both long positions does less to diversify the risk.

For sales-only trading that finishes at zero inventory, the full objective
also contains a fixed offset: $J_{\rm full}=K(\mathbf X)+J$.
It changes the reported total cost, but cannot change which schedule is best.

\textbf{Where we are:} this continuum limit is checked. We optimize $J$ below.

\section*{Step 3 --- Find an equation for the best remaining cost}

Define $\mathcal V(\mathbf x,t)$ as the \textbf{smallest possible remaining
cost} if we hold $\mathbf x$ at time $t$. It is a function of inventory and
time, rather than a functional of a whole path.

Bellman's idea is simple: make the best choice over the next short interval,
then follow the best schedule from the state reached. Expanding this statement
to first order in time gives the Hamilton--Jacobi--Bellman (HJB) equation:
\[
 \boxed{\partial_t\mathcal V
 -\tfrac14(\nabla\mathcal V)^\top H^{-1}\nabla\mathcal V
 +\lambda\mathbf x^\top C\mathbf x=0.}
\]
Here $\nabla\mathcal V$ means derivatives with respect to inventory:
it measures how much extra cost comes from holding a little more of each
asset. The equation is nonlinear because this gradient is squared.

Because our cost is quadratic and our inventory dynamics are linear, we
can solve it with a quadratic value function:
\[
 \mathcal V(\mathbf x,t)=\tfrac12\mathbf x^\top P(t)\mathbf x,
 \qquad
 \boxed{\dot P=\tfrac12PH^{-1}P-2\lambda C,\quad
 \mathbf v^*=\tfrac12H^{-1}P\mathbf x.}
\]
$P(t)$ describes the curvature of the cost surface. Its equation is called a
\emph{matrix Riccati equation}. We have reduced a partial differential
equation (PDE) in inventory and time to an ordinary differential equation
(ODE) for matrix entries that depend only on time.

To set the terminal condition, we can charge a penalty
$\tfrac12\mathbf x(T)^\top\Phi\mathbf x(T)$ for leftover inventory; then
$P(T)=\Phi$. Finite $\Phi$ can leave stock unsold. Exact liquidation is a
separate boundary condition, recovered as the penalty becomes arbitrarily
large in every direction.

\textbf{Where we are:} the HJB and Riccati equations are verified~\cite{bdg}.
They allow buying as well as selling, so we must check that our chosen
solution actually sells throughout. This is the purpose of the next stage.

\newpage
\section*{Stage 3 --- Choose a small example we can check completely}

\emph{Numbering note: ``Step 3'' above is the derivation; ``Stage 3'' in
WorkPlan is the numerical example. They are different tasks.}

We choose a synthetic two-asset example:
\[
 H=\begin{pmatrix}1&0\\0&1\end{pmatrix},\qquad
 C=\begin{pmatrix}1&0.4\\0.4&1\end{pmatrix},\qquad
 \lambda=T=1,\qquad \mathbf X=\begin{pmatrix}3\\1\end{pmatrix}.
\]
These are convenient dimensionless parameters, not measurements from a real
market. The assets have equal individual volatility and trading impact, but
positively correlated price changes. The unequal starting positions make
both collective modes matter.

\subsection*{Think of coupled oscillators and normal modes}

For this example, change coordinates to
\[
 y_+=\frac{x_1+x_2}{\sqrt2},\qquad
 y_-=\frac{x_1-x_2}{\sqrt2}.
\]
These are the symmetric and antisymmetric combinations familiar from two
coupled oscillators. They diagonalize $C$: its eigenvalues are
$\mu_+=1.4$ and $\mu_-=0.6$. Each mode then evolves independently.

With exact liquidation, the solution is
\[
 y_\pm(t)=y_\pm(0)
 \frac{\sinh[\kappa_\pm(T-t)]}{\sinh(\kappa_\pm T)},\qquad
 \kappa_\pm=\sqrt{\lambda\mu_\pm},
\]
\[
 \kappa_+\approx1.1832,\qquad \kappa_-\approx0.7746.
\]
The $\kappa$'s are inverse time scales. Unlike a freely vibrating oscillator,
these inventory modes do not oscillate: we choose paths that end at zero.
The more strongly risk-penalized symmetric mode is reduced more quickly.

\subsection*{Does the solution ever buy stock back?}

We differentiate the inventory path and check
$\mathbf v=-\dot{\mathbf x}$. The all-time proof gives
\[
 \min_{0\le t\le1}v_1(t)\approx2.5063>0,\qquad
 \min_{0\le t\le1}v_2(t)\approx0.6931>0.
\]
Both stocks are therefore sold throughout the interval. For this particular
family of equal-impact, equal-volatility examples with $X_1\ge X_2>0$, the starting ratio must satisfy
\[
 \frac{X_1}{X_2}\le15.9892.
\]
Our ratio is $3$, so it passes comfortably. The proof is in WorkPlan Stage~3.

The optimized running cost is $J^*\approx13.810512$. As an independent
check, we minimized the original discrete quadratic cost on increasingly
fine time grids. At $N=1024$, the largest inventory difference from the
continuous solution is about $1.7\times10^{-8}$.

\textbf{Where we are:} this classical example is verified, including the
sales-only condition for the entire path. It is our reference answer for
testing the later numerical method. This does not establish sales-only
behavior for every possible inventory on a PDE grid.

\newpage
\section*{Stage 4 --- Turn the nonlinear equation into a linear one}

The quantum simulation methods we want to use evolve linear equations.
The squared-gradient term in the HJB is therefore an obstacle.

First use the time remaining, $s=T-t$, so that the terminal data become
initial data at $s=0$. Then introduce a mathematical diffusion coefficient
$\nu>0$. With $W(\mathbf x,s)$ denoting the resulting regularized value
function, the equation is
\[
 \partial_sW=\nu\nabla\!\cdot(H^{-1}\nabla W)
 -\tfrac14(\nabla W)^\top H^{-1}\nabla W
 +\lambda\mathbf x^\top C\mathbf x.
\]
The diffusion acts in \emph{inventory space}, not ordinary physical space.
$\nu$ is an adjustable parameter in the mathematical reformulation; it is
not the stock-price volatility contained in $C$.

Now make the Cole--Hopf change of variables~\cite{kappen}:
\[
 \boxed{\psi=e^{-W/(4\nu)},\qquad \partial_s\psi=-\hat H\psi,}
 \qquad
 \hat H=-\nu\nabla\!\cdot(H^{-1}\nabla)
       +\frac{\lambda}{4\nu}\mathbf x^\top C\mathbf x.
\]
Differentiating the exponential produces a squared-gradient term that
cancels the nonlinear one. The resulting equation is linear in $\psi$.
The matrix $H$ describes trading impact; the operator $\hat H$ acts on
functions of inventory. They are different objects.

Adding diffusion usually changes an optimization problem. Here, on the full
inventory space with fixed finite terminal penalty, it only adds an
inventory-independent term to the quadratic value function. Its gradient
is unchanged. We can therefore recover the same optimal control through
\[
 \boxed{\mathbf v^*=-2\nu H^{-1}\nabla\log\psi.}
\]
This exact statement uses our quadratic structure. It does not automatically
extend to a bounded numerical grid or to a globally sales-constrained problem.

\textbf{Where we are:} the transformation and recovered control are checked
algebraically. Solving the transformed PDE numerically is still to do.

\section*{Stage 5 --- Why a harmonic oscillator appears}

In the normal-mode coordinates of our example,
\[
 \hat H=\sum_{i=+,-}\left[-\nu\frac{\partial^2}{\partial y_i^2}
                    +\frac{\kappa_i^2}{4\nu}y_i^2\right].
\]
This has the familiar harmonic-oscillator form: a second-derivative term
plus a quadratic potential. Its eigenvalues and natural Gaussian widths are
\[
 E_{n_+,n_-}=\kappa_+(n_++\tfrac12)+\kappa_-(n_-+\tfrac12),
 \qquad \ell_i=\sqrt{\frac{2\nu}{\kappa_i}}.
\]
Here the $E$'s describe decay rates in $\psi_s=-\hat H\psi$. This is
\emph{imaginary-time} evolution: a mode decays as $e^{-E s}$ instead of
acquiring a phase $e^{-iEt/\hbar}$ as in ordinary quantum mechanics.

\textbf{Next step:} choose $\nu$ and a sufficiently large, fine inventory
grid to represent these Gaussians and recover their gradients accurately.

\newpage
\section*{Stages 6--9 --- From a PDE to a quantum resource estimate}

\textbf{6. Put inventory space on a grid.}
Represent $\psi(\mathbf x)$ at finitely many inventory values and replace
derivatives by finite differences. This produces a matrix $\hat H_d$.
The grid is in inventory, unlike the time grid used for the original trade
list. Check the recovered selling rate against our classical answer while
refining the grid and enlarging the domain.

\textbf{7. Embed the decay into unitary evolution.}
A quantum state evolves by a norm-preserving unitary transformation, whereas
$e^{-s\hat H_d}$ shrinks norms. Schr\"odingerisation~\cite{jin} adds an
auxiliary coordinate, then uses a Fourier transform to build an enlarged
unitary evolution. Recovering the desired solution from this larger system
is part of the algorithm and has a cost.

\textbf{8. Count the resources.}
If each of the $M$ assets uses $2^{n_x}$ inventory-grid points, the auxiliary
coordinate uses $n_p$ qubits, and the circuit needs $n_{\rm anc}$ working
qubits, then
\[
 N_q=Mn_x+n_p+n_{\rm anc}.
\]
A qubit is a quantum two-state system. We also want $N_T$, the number of
T gates, which are costly operations in fault-tolerant quantum computation.
The subscript here names a gate; it is not the liquidation deadline $T$.
An actual gate count needs an explicit circuit and a target accuracy.

\textbf{9. Read out a useful answer.}
Preparing a quantum state is not the same as reading its whole function.
The encoded amplitudes are normalized:
\[
 a_j=\frac{\psi(\mathbf x_j)}{\sqrt{\sum_k|\psi(\mathbf x_k)|^2}},
 \qquad \Pr(\text{grid point }j)=|a_j|^2.
\]
We want the logarithmic gradient near the starting inventory, because it
gives the selling rate. Estimating ratios of nearby amplitudes may be
expensive when their probabilities are small. This measurement cost must
be included in the total resource estimate.

\section*{Where to resume the project}

The classical solution and the chosen two-asset example are checked.
The Cole--Hopf algebra is checked. We do not yet have a validated numerical
PDE solver or numerical quantum resource counts.

The next concrete calculation is to \textbf{choose finite terminal data,
solve the linear PDE on a small classical grid, and recover the selling
rate}. Compare it with the Riccati solution using the \emph{same} terminal
penalty. A finite penalty must not be compared as if it imposed exact
liquidation. Then measure how the error changes with grid spacing, domain
size, and $\nu$.

To rerun the existing checks in this checkout:
\texttt{.venv/bin/python -B verify\_research.py}.
Full derivations, corrected finite-penalty formulas, source checks, and
remaining assumptions are in \texttt{WorkPlan.tex}.

\begingroup
\small
\begin{thebibliography}{9}
\bibitem{ac} R.~Almgren and N.~Chriss, \emph{Optimal Execution of Portfolio
Transactions} (2000). Trading objective and classical trajectories.
\bibitem{bdg} P.~Bergault, F.~Drissi and O.~Gu\'eant (2022).
Multi-asset HJB and Riccati equations; Appendix A is the relevant special case.
\url{https://arxiv.org/abs/2103.13773}.
\bibitem{kappen} H.~J.~Kappen (2005). Exponential transformation of suitable
stochastic-control equations. \url{https://arxiv.org/abs/physics/0411119}.
\bibitem{jin} S.~Jin, N.~Liu and Y.~Yu (2023).
Schr\"odingerisation: technical details, especially Section 3.1.
\url{https://arxiv.org/abs/2212.14703}.
\end{thebibliography}
\endgroup
\end{document}
